\documentclass[de]{../../elegantbook}
% ========== Titelblatt und Grundinformationen ==========
\title{ElegantBook Testdokument}
\subtitle{Kernmerkmale der Vorlage}
\author{Testingenieur}
\institute{Vorlagen-Testgruppe}
\date{\today}
\version{v1.0}
\extrainfo{Dieses Dokument dient ausschließlich zum Testen von Vorlagenfunktionen und hat keinen akademischen Wert}
\bioinfo{Testcharge}{TEST-FULL-20260506}
\logo{../../figure/logo-blue.png}
\cover{../../figure/cover.jpg}

% ========== Inhaltsverzeichnis und Grundeinstellungen ==========
\setcounter{tocdepth}{3} % Inhaltsverzeichnis bis Ebene 3 anzeigen

% ========== Test-Literaturverzeichnis ==========
\addbibresource[location=local]{../en/reference.bib}

\begin{document}

\maketitle
\frontmatter
\tableofcontents

% ========== Vorwort (mit Abschnittsgliederung) ==========
\pagenumbering{Roman}
\setcounter{page}{0}
\chapter*{Testvorwort}
\markboth{Testvorwort}{}
\section*{Dokumentbeschreibung}
\markright{Dokumentbeschreibung}
Dieses Dokument ist ein vollständiger Testfall für ElegantBook und deckt alle grundlegenden und erweiterten Funktionen der Vorlage ab. Alle Inhalte sind fiktive Testdaten ohne praktische Bedeutung.
\newpage
\section*{Testumfang}
\markright{Testumfang}
Der Test umfasst: Titelblatt, Inhaltsverzeichnis, Schriftarten, mehrstufige Listen, komplexe mathematische Formeln, alle mathematischen Umgebungen, Querverweise, Abbildungen und Tabellen, Randbemerkungen, Kapitelzusammenfassungen, Übungen, Literaturverzeichnis und Anhänge.

\mainmatter

% ========== Kapitel 1: Grundsatzlayout und mehrstufige Listen ==========
\chapter{Grundlayout- und Listentest}
\begin{introduction}[Testpunkte dieses Kapitels]
\item Chinesische Schriftarten und Textformatierung
\item Verschachtelung ungeordneter Listen mit 3 Ebenen
\item Verschachtelung geordneter Listen mit 3 Ebenen
\item Normale Absätze und hervorgehobener Text
\end{introduction}

\section{Schriftarten- und Textest}
Normaler Fließtext-Test. \textbf{Fett}, \textit{Kursiv}, \underline{Unterstrichen}, \textcolor{red}{Roter Text-Test}.

\section{Test ungeordnete Liste mit 3 Ebenen}
\begin{itemize}
\item Ebene-1-Testpunkt A
\item Ebene-1-Testpunkt B
    \begin{itemize}
    \item Ebene-2-Testpunkt B1
    \item Ebene-2-Testpunkt B2
        \begin{itemize}
        \item Ebene-3-Testpunkt B2-1
        \item Ebene-3-Testpunkt B2-2
            \begin{itemize}
            \item Ebene-4-Testpunkt B2-1a
            \item Ebene-4-Testpunkt B2-2b
            \end{itemize}
        \end{itemize}
    \end{itemize}
\end{itemize}

\section{Test geordnete Liste mit 3 Ebenen}
\begin{enumerate}
\item Ebene-1-geordneter Test 1
\item Ebene-1-geordneter Test 2
    \begin{enumerate}
    \item Ebene-2-geordneter Test 2-1
    \item Ebene-2-geordneter Test 2-2
        \begin{enumerate}
        \item Ebene-3-geordneter Test 2-2-1
        \item Ebene-3-geordneter Test 2-2-2
            \begin{enumerate}
            \item Ebene-4-geordneter Test 2-2-2-1
            \item Ebene-4-geordneter Test 2-2-2-2
            \end{enumerate}
        \end{enumerate}
    \end{enumerate}
\end{enumerate}

% ========== Kapitel 2: Test komplexer mathematischer Formeln ==========
\chapter{Test komplexer mathematischer Formeln}
\begin{introduction}
\item Einzeilige komplexe Formel
\item Mehrzeilige Formeln (align)
\item Matrizen / Determinanten
\item Mehrfachintegrale / Summen / Grenzwerte / Brüche
\item Formelbeschriftungen und Querverweise
\end{introduction}

\section{Einzeilige komplexe Formel}
Inline komplexe Formeln: $\sum_{i=1}^n \frac{x_i^2}{\sqrt{y_i+1}} \geq \bar{x}^2$, $\lim_{n\to\infty} \left(1+\frac{1}{n}\right)^n = e$.

\section{Mehrzeilige Formeln und Matrizen}
\begin{align}
a^2 + b^2 &= c^2 \label{eq:gougu} \\
\int_{0}^{1}\int_{0}^{1} f(x,y) dxdy &= 1 \label{eq:double} \\
\begin{vmatrix}
a & b \\
c & d
\end{vmatrix} &= ad-bc \label{eq:det}
\end{align}

\section{Formeln mit großen Operatoren}
\begin{equation}\label{eq:series}
\sum_{k=0}^{\infty} \frac{x^k}{k!} = e^x,\quad \oint_{C} Pdx+Qdy = \iint_{D}\left(\frac{\partial Q}{\partial x}-\frac{\partial P}{\partial y}\right)dxdy
\end{equation}

Formel-Querverweis: Satz des Pythagoras \eqref{eq:gougu}, Doppelintegral \eqref{eq:double}, Determinante \eqref{eq:det}, Reihe \eqref{eq:series}.

% ========== Kapitel 3: Vollständiger Test aller mathematischen Umgebungen ==========
\chapter{Vollständiger Test mathematischer Umgebungen}
\begin{introduction}
\item Definition / Satz / Lemma / Korollar / Axiom / Postulat
\item Proposition / Beispiel / Problem / Übung
\item Hinweis / Bemerkung / Schlussfolgerung / Annahme / Eigenschaft
\item Lösung / Beweisumgebung
\end{introduction}

\section{Umgebungstest}

% 1. Definition
\ifdefined\simpleTest
\begin{definition}[Testdefinition]\label{def:test}
\else
\begin{definition}[Testdefinition]{test}
\fi
Dies ist der Testinhalt der Definitionsumgebung ohne tatsächliche mathematische Bedeutung, nur zur Stilprüfung.
\end{definition}

% 2. Satz
\ifdefined\simpleTest
\begin{theorem}[Testsatz]\label{thm:test}
\else
\begin{theorem}[Testsatz]{test}
\fi
Wenn $x>0$, dann gilt $x+1>1$, entsprechend Formel \eqref{eq:gougu}.
\end{theorem}

% 3. Lemma
\ifdefined\simpleTest
\begin{lemma}[Testlemma]\label{lem:test}
\else
\begin{lemma}[Testlemma]{test}
\fi
Für jede reelle Zahl $x$ gilt $x^2\geq0$.
\end{lemma}

% 4. Korollar
\ifdefined\simpleTest
\begin{corollary}[Testkorollar]\label{cor:test}
\else
\begin{corollary}[Testkorollar]{test}
\fi
Aus Lemma \ref{lem:test} folgt: $(x-1)^2\geq0$.
\end{corollary}

% 5. Axiom
\ifdefined\simpleTest
\begin{axiom}[Testaxiom]\label{axi:test}
\else
\begin{axiom}[Testaxiom]{test}
\fi
Testaxiom ohne praktische Bedeutung.
\end{axiom}

% 6. Postulat
\ifdefined\simpleTest
\begin{postulate}[Testpostulat]\label{pos:test}
\else
\begin{postulate}[Testpostulat]{test}
\fi
Testpostulat nur zur Umgebungsprüfung.
\end{postulate}

% 7. Proposition
\ifdefined\simpleTest
\begin{proposition}[Testproposition]\label{pro:test}
\else
\begin{proposition}[Testproposition]{test}
\fi
Inhalt der Testproposition, verknüpft mit Definition \ref{def:test}.
\end{proposition}

% 8. Beispiel
\begin{example}\label{exa:test}
Beispielumgebungstest mit automatischer Nummerierung.
\end{example}

% 9. Problem
\begin{problem}\label{probl:test}
Problemungebungstest zur Stellung ungelöster Aufgaben.
\end{problem}

% 10. Übung
\begin{exercise}\label{exer:test}
Berechnen Sie: $1+2+3+\dots+10$.
\end{exercise}

% 11. Hinweis
\begin{note}
Hinweisumgebung ohne Nummerierung für wichtige Erinnerungen.
\end{note}

% 12. Bemerkung
\begin{remark}
Bemerkungsumgebung für ergänzende Erklärungen.
\end{remark}

% 13. Schlussfolgerung
\begin{conclusion}
Schlussfolgerungsumgebung für zusammenfassenden Testinhalt.
\end{conclusion}

% 14. Annahme
\begin{assumption}
Annahmeumgebung zur Festlegung von Testbedingungen.
\end{assumption}

% 15. Eigenschaft
\begin{property}
Eigenschaftsumgebung zur Beschreibung von Testmerkmalen.
\end{property}

% 16. Lösung
\begin{solution}
Lösung zu Übung \ref{exer:test}: Die Summe beträgt $55$ (Testantwort).
\end{solution}

% 17. Beweis
\begin{proof}
Beweis zu Satz \ref{thm:test}: Da $x>0$ ist, folgt durch Addition von 1 auf beide Seiten die Behauptung.
\end{proof}

% ========== Kapitel 4: Querverweise · Abbildungen · Randbemerkungen ==========
\chapter{Querverweise und erweiterte Funktionen}
\begin{introduction}
\item Alle Arten von Querverweisen
\item Einfacher Abbildungs- und Tabellentest
\item Randbemerkungsfunktionstest
\item Kapitelnavigationsprüfung
\end{introduction}

\section{Zusammenfassung aller Querverweise}
Definition \ref{def:test}, Satz \ref{thm:test}, Lemma \ref{lem:test}, Korollar \ref{cor:test}, Axiom \ref{axi:test}, Postulat \ref{pos:test}, Proposition \ref{pro:test}, Beispiel \ref{exa:test}, Problem \ref{probl:test}, Übung \ref{exer:test}.

\section{Abbildungs- und Tabellentest}
\begin{figure}[h]
    \centering
    \includegraphics[width=0.5\textwidth]{../../image/scatter.jpg}
    \caption{Testabbildung}\label{fig:test}
\end{figure}
\begin{table}[h]
    \centering
    \begin{tabular}{c*{6}{c}}
    \toprule
        $n$ & 1 & 2 & 3 & 4 & 5 & 6 \\
    \midrule
        $a_n$ & 1 & 1 & 2 & 3 & 5 & 8 \\ 
    \bottomrule
    \end{tabular}
    \caption{Testtabelle}\label{tab:test}
\end{table}
Abbildungs-/Tabellenverweis: \figref{fig:test} ist die Testabbildung, \tabref{tab:test} ist die Testtabelle.

% ========== Kapitelabschlussübungen ==========
\begin{problemset}[Kapitel-Testübungen]
\item Überprüfen Sie das Anzeigestil der Definition \ref{def:test}
\item Herleiten Sie die Formel \eqref{eq:double}
\item Prüfen Sie, ob Abbildung \ref{fig:test} korrekt angezeigt wird
\end{problemset}

% ========== Literaturverzeichnis ==========
\nocite{*}
\printbibliography

% ========== Anhang ==========
\appendix
\chapter{Testanhang}
Der Anhanginhalt dient zur Prüfung von Anhangskapiteln, Kopfzeilen, Seitenzahlen und Layoutstil ohne tatsächliche Informationen.
\section{Testanhang 1}
Der Anhanginhalt dient zur Prüfung von Anhangskapiteln, Kopfzeilen, Seitenzahlen und Layoutstil ohne tatsächliche Informationen.
\newpage
\section{Testanhang 2}
Der Anhanginhalt dient zur Prüfung von Anhangskapiteln, Kopfzeilen, Seitenzahlen und Layoutstil ohne tatsächliche Informationen.
\end{document}